%!TEX root = main.tex

\section{Our adaptive-based broadcast algorithm for heterogeneous MANET's}

\subsection{Adaptation mechanism}

\raziel{
Lemma to prove the correctness of our adaptation mechanism. Plan of work:
\begin{itemize}
	\item write down lemma having broadcast messages as abstract as possible. Some consideration while writing lemma:
	\begin{itemize}
		\item the generalization of our claim will take place saying that our mechanism is valid when $A$ is the set of ALL existing broadcast algorithms and not when $A$ is just ONE set of protocols, as we do right now.
		%
		This can be mention as future work having a reasoning about what exactly is $A$ in this current version of the lemma.
		%
		 Q: How can we proof a lemma when $A$ refers to all existing broadcast protocols? (big question)
		\begin{itemize}
			\item Q: shall we have a taxonomy of broadcast algorithms in mind? ANSWER: say that we do.
			%
			Then, we must define that any $p \in O$ or $p \in N$ where $O$ refers to all overlay-based algorithms and $N$ is the set of non overlay-based algorithms; which is a valid generalization(I guess?)
			\item proof by induction over any $p_{i}$
			\item proof by contradiction over any $p_{i}$. $\exists p_{i} \in O$ that lost a message $M_{p_{j}}$, then find a bimodal value in data structures of $p_{i}$ OR an incorrect functioning in the way $p_{i}$ operates to show that a contradiction actually exists
		\end{itemize}
	\end{itemize}
	\item write down proof of lemma by contradiction
\end{itemize}
}

\begin{lemma}
	Let be $M_{p_{i}}$ a message and $p_{i}$ a broadcast protocol such that $p_{i} \in A$, where $A$ is a set of protocols. If a node $N$ receives $M_{p_{i}}$ while $p_{j}$ is in execution at $N$, then Algorithm~\ref{alg:adaptation_mechanism} translates $M_{p_{i}}$ into $M_{p_{j}}$ for being disseminated by $N$ using $p_{j}$.
\label{lemma:adapt_mech}
\end{lemma}

We need to guarantee that Algorithm~\ref{alg:adaptation_mechanism} does not drop any incoming message $M_{p_{i}}$ while a protocol $p_{j}$ is in execution on node $N$.
%
This potential loss of messages could take place only in line~\ref{alg:adaptation_mechanism:handle} or ~\ref{alg:adaptation_mechanism:emit} of Algorithm~\ref{alg:adaptation_mechanism}.
%
We avoid the case where there is a call of function $\AlgoHandle(\Ms)$ (line~\ref{alg:adaptation_mechanism:handle}, Algorithm~\ref{alg:adaptation_mechanism}), because our adaptation mechanism do not interfere with the way any protocol operates.
%
Thus, we offer a proof of lemma~\ref{lemma:adapt_mech} considering the case where the message translation take place (line~\ref{alg:adaptation_mechanism:emit}, Algorithm~\ref{alg:adaptation_mechanism}), as follows:

\textbf{Proof.}
Let be $M_{p_{i}}$ a message that was successfully received by node $N$ for the first time, but its translation into message $M_{p_{j}}$ was not forwarded.
%
Given that we want to validate the message translation procedure, not unexpectedly, protocol $p_{j}$ is in execution on node $N$. 
%
We follow that the condition in line~\ref{alg:adaptation_mechanism:eval_ids} of Algorithm~\ref{alg:adaptation_mechanism} is not met because $p_{i} \neq p_{j}$ and a timer $T$ with duration \AdaptTimer starts (line~\ref{alg:adaptation_mechanism:set_timer}).
%
Eventually, $T$ will expire to obtain $M_{p_{j}}$ from message $M_{p_{i}}$, a call of function $\AlgoEmit()$ have place and $M_{p_{j}}$ is forwarded (line~\ref{alg:adaptation_mechanism:emit}).
%
We then reach a contradiction, because our first assumption was that the retransmission of $M_{p_{j}}$ didn't take place. We then conclude that such a message $M_{p_{j}}$ does not exist$\qed$.

In Section~\ref{} we analyze to which extent the policies of retransmission from broadcast protocols have an impact in network coverage.


\subsection{Cost of control messages}
Control messages serve to fetch fresh information about nodes' surroundings with the aim of deciding when to forward broadcast messages.
%
Control messages are critical in the latest adaptive-based techniques to broadcast in homogeneous MANET's ~\cite{hypergossiping, ProbFlood} \raziel{If you like the idea, add more referece to adaptive-based approaches, otherwise, to contradict this idea feel free to reference works where adaptation takes place without the use of control messages}.
%
However, these techniques perform a pure flooding approach in a periodic way to disseminate control messages.
%
As a result, \emph{storms of broadcast and control messages} reduces network coverage \raziel{I propose to rename "network coverage" as "reachability", the latter has been used by the community in seminal works ~\cite{hypergossiping, BroadcastStormProblem}. Reachability will give us the proportion of nodes that successfully received a broadcast message in all partitions of our experimental scenario. Or, we can keep add "reachability" to our list of broadcast metrics.}

In \OurAdaptiveTech, we propose the use of a beacon-less technique (\raziel{An augmented version of ABBA where timers are inversaly proportinal to nodes' density. Particulary, we are interested in maintaning the list of neighbors on each node. Additionally, the upperbounds values to set timers can be dinamically changed too for speeding-up the dissemination time; this is important when a switch of context take place.}) to disseminate control messages with the aim of having updated lists of nodes' neighbors. We formulate the next hypothesis: \textbf{the overall energy consumption of nodes used by a the periodic simple-flooding technique to disseminate control messages is bigger than the cost of control messages dissemination with ABBA $+$ $f_{GPS}(N)$}, where $f_{GPS}$ approximates the cost of fetching $N$' GPS coordinates periodically. This frequency is proportional to $N$ velocity (\raziel{we can gatter this information from mobile's accellerometer}). This will require a deep analysis in ABBA (\raziel{in form a table to map the inpact in our broadcast metrics and the different tunes of ABBA}).

\paragraph{Drawbacks} ABBA requires nodes position to its functioning. 

\subsection{Brainstorming}
\paragraph{Ideas to ameliorate cost of localization}
This work~\cite{enloc} ameliorates the energy cost of getting localization of nodes.

Current adaptive solutions change the value of thresholds to be well informed before performing message retransmissions.

\paragraph{Some questions to consider in adaptation:}
\begin{itemize}
	\item How to get a fresh list of nodes' neighbors?
	\item When to perform a switch of context (change of algorithm)?
	\item How to disseminate SWITCH\_OF\_CONTEXT messages?
\end{itemize}
